\makeabstract
    {
    MYC Super Enhancer Activity in in T-ALL Requires TCF-1
    }
    {
    Benjamin Edwards\textsuperscript{1,2}, Vasili Toliopoulos\textsuperscript{1}, George Tousinas\textsuperscript{1}, Fotini Gounari\textsuperscript{1}
    }
    {
    \textsuperscript{1} Department of Biology, Amherst College, Amherst, MA\\
\textsuperscript{2} Department of Immunology, Mayo Clinic, Phoenix, AZ
    }
    {
    Aberrant activation of MYC is a common oncogenic event across human cancers, including T-cell acute lymphoblastic leukemia (T-ALL). In approximately 5\% of human T-ALL cases, MYC overexpression results from duplication and activation of N-ME, a distal Myc super-enhancer. We have developed a spontaneous mouse model of T-ALL driven by T-cell-specific stabilization of \ensuremath{\beta}-catenin combined with the loss of the transcription factor HEB (CATHEB\ensuremath{\Delta}). This model faithfully recapitulates spontaneous N-ME super enhancer duplication and MYC overexpression during the early double positive (DP blast) stage of thymocyte development. We have also identified novel recruitment of TCF-1 to the N-ME super enhancer, and genetic ablation of TCF-1 from CATHEB\ensuremath{\Delta} thymocytes completely prevents leukemogenesis. However, it remained unclear whether this protection resulted from disruption of thymocytes development, or whether TCF-1 has a key role in N-ME superenhancer activation via chromatin accessibility and remodeling.
    }
    {
    To distinguish between these possibilities we characterized thymocytes development in CATHEB\ensuremath{\Delta}TCF1\ensuremath{\Delta} mice using multiparameter flow cytometry with CD4, CD8, CD69, CCR7, CD25, TCR\ensuremath{\beta}, TCF-1, and viability staining. 
    }
    {
    Although total thymocytes numbers were significantly reduced consistent with the established role of TCF-1 in thymocyte development, CATHEB\ensuremath{\Delta}TCF1\ensuremath{\Delta} thymocytes progressed to the DP blast stage that is targeted for leukemogenesis in CATHEB\ensuremath{\Delta} mice. In addition, these mice exhibited significantly increased proportions of mature, post-DP T cells expressing CD69+ and CCR7+.
    }
    {
    These findings demonstrate that loss of leukemogenesis in CATHEB\ensuremath{\Delta}TCF1\ensuremath{\Delta} mice cannot be explained solely by a developmental block at the DP stage. Instead, they support a direct role for TCF-1 in promoting leukemic transformation, consistent with its proposed function in activating the N-Me super-enhancer and driving MYC-dependent T-ALL.
    }

    \newpage
    